\section{磁单极子的有效超势}\label{supermon}

\begin{keytakeaways}
\begin{itemize}\setlength{\itemsep}{2pt}
  \item \textbf{物理图像}：$\R^2\times\R_+$ 上的光滑 BPS 磁单极子与边界上的 Nahm 极点相互作用时不存在精确解；Nahm 极点把磁单极子推向 $y\to\infty$，使其在指数精度上“几乎是解”，故其低能行为由一个有效超势支配。
  \item \textbf{数学工具}：用一对手征超场 $\sW$（$\R^2$ 中的位置）与 $\sY=\ca_1y+i\vartheta$（$y$ 位置与 $U(1)$ 集体坐标的组合）参数化单磁单极子模空间；用多项式 $K(z)=\prod_a(z-z_a)^{k_a}$ 记录边界奇异磁单极子的磁荷；通过积掉大质量场 $\sY$ 生成 $\sW$ 的有效超势。
  \item \textbf{主要结论}：有效超势为 $\W=\ca\sW+\zeta\sY+\Lambda K(\sW)\exp(-\sY)$；其临界点条件逐一复现第 2、3 节关于解存在性的定性结论；积掉 $\sY$ 后得到 Yang-Yang 函数 (\ref{longyang})，即 $\ca\sum_i\sW_i+\zeta\sum_{i,a}\log(\sW_i-z_a)-2\zeta\sum_{i<j}\log(\sW_i-\sW_j)$。
  \item \textbf{技术要点}：多个光滑磁单极子的 $\sW_i$ 重合时，$\sY_i$ 不再是好坐标，必须改用由散射矩阵定义的 $\nY_i$；两者之间的变换 (\ref{rela}) 在 $\sW_i=\sW_j$ 时发散。
\end{itemize}
\end{keytakeaways}

\subsection{结果概述}\label{outline}

从第~\ref{symbr}~节开始，我们把三维超对称方程的解解释为包含光滑 BPS 磁单极子的位形。然而，那些讨论纯粹是定性的。本节将把上述推理定量化。我们将为 $\R^2\times \R_+$ 上与 Nahm 极点以及边界上的奇异磁单极子相互作用的光滑 BPS 磁单极子构造一个有效超势。这个有效超势将直接解释第 2 节与第 3 节中关于在参数取各种值时我们的方程有解或无解的全部定性结论。此外，通过积掉一些大质量场，我们将能够还原一直是主要工具之一的 Yang-Yang 函数 (\ref{yy})。

我们先考虑作为出发点的超对称方程
\begin{align}\label{bpstwo} (F-\phi\wedge \phi + \t \,\d_A\phi)^+&=0   \cr
                             (F-\phi\wedge\phi-\t^{-1}\d_A\phi  )^- &=0 \cr
                               \d_A\star\phi&=0,                     \end{align}
它定义在 $\R\times \R^3$ 上（其中第一个因子是“时间”方向），并问光滑磁单极子的 Bogomolny 方程的解如何能嵌入为这些方程的解。当且仅当 $\t=1$（或 $-1$）时这才可能：此时 $\phi$ 的唯一非零分量是时间分量，我们这里记作 $\phi_t$（而不像前面那样记作 $\phi_1$），同时把 $A$ 的时间分量取为零。于是方程 (\ref{bpstwo}) 约化为 Bogomolny 方程
\begin{equation}\label{bog} F=\star \d_A\phi_t.\end{equation}

要使 Bogomolny 方程有光滑磁单极子解，场 $\phi_t$ 在无穷远处必须具有期望值，这意味着第~\ref{symbr}~节的实对称破缺参数 $\ca_1$ 必须非零。另一方面，第~\ref{simvar}~节的复对称破缺参数 $\ca$ 必须为零（否则 $\phi_t$ 就不会是 $\phi$ 的唯一非零分量）。


$\R^3$ 上 Bogomolny 方程的基本解是 $G=SO(3)$ 或 $SU(2)$ 的单磁单极子解。它的磁荷为 $\m=2$（以基本奇异磁单极子的磁荷为 $\m=1$ 的单位制）。单磁单极子解的模空间是 $\P=\R^3\times S^1$，其中 $\R^3$ 刻画磁单极子在空间流形 $\R^3$ 上的质心位置，$S^1$ 由无穷远处未破缺的 $U(1)$ 规范对称性的集体坐标 $\vartheta$ 参数化。当然，$\P$ 是一个 hyper-Kähler 流形；这可由 $\R^3$ 上 Bogomolny 方程的未破缺超对称推出。不过，由于我们将考虑破坏部分超对称的微扰，就本文目的而言，更可取的做法是仅把 $\P$ 看作其某个复结构下的复流形。我们把 $\R^3$ 分解为 $\R^2\times \R$（最终将用 $\R^2\times \R_+$ 代替它）。光滑磁单极子沿 $\R^2$ 的运动由一个手征超场 $\sW$ 参数化。而光滑磁单极子在 $\R$ 方向的位置 $y$ 与集体坐标 $\vartheta$ 组合成第二个手征超场
\begin{equation}\label{work}\sY=\ca_1 y+i\vartheta.   \end{equation}

现在我们要构造描述这一情形的有效超势 $\W(\sW,\sY)$。若 $\ca=\zeta=0$，则对 $\sW$ 与 $\sY$ 的每个取值都存在光滑磁单极子解，这意味着 $\sW$ 与 $\sY$ 的每个取值都是 $\W$ 的临界点，因此 $\W$ 必须为零（模去一个无关常数）。反之，只要 $\ca$ 或 $\zeta$ 非零，就不存在超对称磁单极子解，即 $\W$ 没有临界点。但引入 $\ca$ 与 $\zeta$ 并不破坏“给 $\sW$ 或 $\sY$ 加常数”的对称性。因此 $\W$ 在 $\sW$ 或 $\sY$ 作常数平移下必须只改变一个加性常数（即模加性常数不变）；换言之，$\W$ 必须是 $\sW$ 与 $\sY$ 的线性函数。事实上，$\W$ 的形式是
\begin{equation}\label{thelf}\W=\ca \sW+\zeta \sY.    \end{equation}
这可由如下考虑推出。$\R^2$ 转动下的不变性要求 $\ca$ 与 $\sW$ 只能以乘积 $\ca \sW$ 的形式出现。另一方面，虽然在 $\zeta=0$ 时超势 $\W$ 在微观上由单值函数 (\ref{superone}) 给出，但一旦 $\zeta$ 非零，它就由 Chern-Simons 函数 (\ref{supertwo}) 给出，而后者只在模 $2\pi i\Z$ 的意义下单值。由于 $\sY$ 同样只在模 $2\pi i\Z$ 的意义下单值（因为 $\vartheta$ 是角度变量），有效超势 (\ref{thelf}) 的多值性与微观描述的多值性恰好一致。

现在我们令 $\ca$ 与 $\zeta$ 为零，但把 $\R^2\times \R$ 换成 $\R^2\times \R_+$，其中 $\R_+$ 是半直线 $y\geq 0$，并假定 $y=0$ 处有通常的 Nahm 极点。我们假定 $y\to\infty$ 时发生对称破缺，$\phi_t\to {\mathrm{diag}}(\ca_1,-\ca_1)$，而 $A$ 以及（在 $\ca=0$ 时）$\vec\phi$ 当然在 $y\to\infty$ 时趋于零。在任何解中，各场在 $y\to\infty$ 时都以指数速度趋近这些渐近值。这反映了如下事实：在 $y=\infty$ 附近，规范对称性被 $\phi_1$ 的期望值从 $SU(2)$ 约化为 $U(1)$，所有带电场的质量都正比于 $|\ca_1|$。特别地，在既无奇异磁单极子也无光滑磁单极子时，完整系统 (\ref{bpstwo}) 的解由 Nahm 方程的一个解给出（相关解见文献~\cite{Krontwo}），其中各场在 $y\to\infty$ 时以指数速度趋近其真空值。

% [Tutor Note: 原文笔误] 此处作者写 $\phi_1$ 的期望值把 $SU(2)$ 破缺到 $U(1)$，但本节开头已把时间分量改记为 $\phi_t$（“rather than $\phi_1$, as before”）。按上下文此处应为 $\phi_t$；原式如此，未改。

接下来，仍取 $\ca=\zeta=0$，我们加入一个位于 $\sW$、$\sY$ 处的光滑磁单极子。光滑磁单极子并不存在精确解，因为 Nahm 极点迫使 $\vec\phi$ 的带电分量取非零值，从而把磁单极子“排斥”到 $y=\infty$。不过，这些带电场在大 $y$ 处指数地小，因此位于大 $y$ 处的光滑磁单极子与解之间的差别是指数小的。既然如此，位于大 $y$ 处的光滑磁单极子必定由某个有效超势支配。该超势必须是 $\sY$ 的单值函数，并在 $y\to\infty$ 时指数地趋于零。指数函数 $\exp(-n\sY)$（$n=1,2,3\dots$）的线性组合满足这些条件。然而，下文很快会看出，预期的定性图像成立当且仅当在 $\ca=\zeta=0$ 时 $\W$ 对 $\exp(-\sY)$ 是线性的：
\begin{equation}\label{forg} \W=\Lambda \exp(-\sY),\end{equation}
其中 $\Lambda$ 为某个常数。

% [Tutor Note: 省略论证] 作者只以“下文很快会看出”带过“为何必须对 $\exp(-\sY)$ 线性”。理由在于：$\W$ 若含 $\exp(-n\sY)$（$n\ge2$）项，则 $\partial\W/\partial\sY=0$ 一般给出多个 $\exp(-\sY)$ 的解，临界点不再由 $\ca=0,\ \zeta\not=0$ 唯一确定，与第~\ref{realsym}~节式 (\ref{onl}) 的单参数族不符；而线性项保证 (\ref{norg}) 对 $\exp(-\sY)$ 有唯一解。

原则上，可以通过计算由大 $y$ 处的光滑磁单极子与 Nahm 极点组成的近似解的微观超势 (\ref{superone}) 来得到这个结果。指数小的项应当来自边界与磁单极子之间的 $W$ 玻色子交换。本文不走这条计算路线，而采取一条捷径：我们将用膜位形来表现光滑磁单极子与 Nahm 极点，在那一框架下，指数型超势 (\ref{forg}) 将来自一个简单的膜瞬子。

我们把这一分析推迟到第~\ref{realm}~节，先讨论 (\ref{forg}) 的含义。首先假定 $\ca=\zeta=0$，于是 (\ref{forg}) 就是完整的超势。立刻可以看出 $\W$ 没有临界点，因此在光滑磁单极子存在时没有超对称解。这与第~\ref{symbr}~节的分析完全一致。在只有实对称破缺且 $\zeta=0$ 时，超对称解被期待是唯一的（即使 $y=0$ 处有奇异磁单极子，这一点仍然成立，而奇异磁单极子尚未纳入 $\W$），并且不要求存在光滑磁单极子。现在来看引入 $\ca$ 与 $\zeta$ 后会发生什么。我们简单地把迄今得到的各项相加，构造完整超势：
\begin{equation}\label{morg}\W=\ca \sW+\zeta\sY +\Lambda\exp(-\sY).\end{equation}
这样引入 $\ca$ 与 $\zeta$ 的理由与前面相同（$\Lambda$ 现在可以依赖于 $\zeta$）。我们来寻找临界点。立刻可以看出，除非 $\ca=0$ 且 $\zeta\not=0$，否则没有临界点。若这些条件满足，则存在一个由 $\sW$ 参数化的单参数临界点族，满足
\begin{equation}\label{norg} \exp(-\sY)=\frac{\zeta}{\Lambda}.\end{equation}
（$\exp(-\sY)$ 的解之所以唯一，正是因为 (\ref{forg}) 中取 $\W$ 对 $\exp(-\sY)$ 线性。）我们在第~\ref{realsym}~节遇到过类似结果，那里在式 (\ref{onl}) 中描述了一族恰好当 $\ca=0$ 且 $\zeta\not=0$ 时存在的解。所涉及的解对应于一个只含实对称破缺的 Miura oper、一个位于任意点 $w\in \R^2$ 的单个 Bethe 根，以及没有奇异磁单极子。我们把 Bethe 根 $w$ 解释为超场 $\sW$ 的取值——换言之，光滑磁单极子在 $\R^2$ 中的位置——而光滑磁单极子在 $\R_+$ 方向的位置由 (\ref{norg}) 确定。这一描述之所以合理，是因为对小的 $\zeta$，光滑磁单极子位于大 $y$ 处，而那里有效超势 $\W$ 是适用的。

下一步是纳入奇异磁单极子。如第~\ref{addsing}~节那样，当复平面上位置为 $z_a,\,a=1,\dots,d$ 处存在磁荷为 $k_a$ 的奇异磁单极子时，引入多项式 $K(z)=\prod_{a=1}^d(z-z_a)^{k_a}$ 是方便的。在第~\ref{scattering}~节我们将论证，奇异磁单极子的唯一效应是把指数项乘以 $K(\sW)$，于是超势变为
\begin{equation}\label{bigsuper}\W=\ca \sW+\zeta\sY +\Lambda K(\sW)\exp(-\sY).\end{equation}
现在我们来考察这一公式的含义。先假定 $\zeta=0$。此时临界点的条件是
\begin{equation}\label{igsuper}K(\sW)=0=\ca+\Lambda K'(\sW)\exp(-\sY).\end{equation}
第一个条件说明 $\sW$ 必须等于某个 $z_a$。当 $\ca=0$ 且所有 $k_a=1$ 时，第二个条件无解。若某个 $a$ 有 $k_a>1$，仍取 $\ca=0$，则第二个条件对任意的 $Y$ 都满足。这一切与我们在第~\ref{solutions}~节和第~\ref{symbr}~节所得一致。现在假定 $\ca\not=0$。若 $k_a=1$，则 (\ref{igsuper}) 中的第二个条件唯一地确定 $Y$；并且若 $\ca$ 小，则 $Y$ 大，从而分析是可靠的。若 $k_a>1$，则 (\ref{igsuper}) 中的第二个条件无法满足。所有这些陈述都与第~\ref{simvar}~节从相当不同的角度所得的结果相符。\footnote{在第~\ref{simvar}~节，我们对一般的 $k_a$ 与 $\ca\not=0$ 找到了一个在 $W=z_a$ 处没有光滑磁单极子的解，以及一个有 $k_a$ 个光滑磁单极子的解。（它们对应于以两种方式解 $PQ=(z-z_a)^{k_a}$，使得 $P$ 与 $Q$ 在 $z=z_a$ 处没有公共零点。）由于 (\ref{bigsuper}) 只是单个光滑磁单极子的超势，所以只有当 $k_a=1$ 时它才描述一个有 $k_a$ 个光滑磁单极子的解。对任意 $k_a$ 的一般分析可以借助任意数目光滑磁单极子的超势 (\ref{giantsuper}) 作出，并与第~\ref{simvar}~节对照。}

% [Tutor Note: 易错点] (\ref{igsuper}) 中第二个条件里的 $K'(\sW)$ 与 $\exp(-\sY)$ 同号相乘；由于 $\Lambda$ 可依赖 $\zeta$ 且其符号未定，$Y$ 究竟取大正还是大负，须由 $\zeta/\Lambda$ 的符号决定，不能想当然地认为“$\ca$ 小 $\Rightarrow$ $Y$ 大正”。

现在考虑 $\zeta\not=0$ 的情形。最有启发性的做法是积掉大质量场 $\sY$，生成 $\sW$ 的有效超势。对固定的 $\sW$，条件 $\partial\W/\partial \sY=0$ 有唯一解 $\exp(-\sY)=\zeta/\Lambda K(\sW)$。把 $\sY$ 取为此值并计算 $\W$，得到（模去一个无关常数）
\begin{equation}\label{twig} \W=\ca \sW + \zeta\log K(\sW)=\ca\sW+\zeta\sum_a k_a
\log(\sW-z_a).  \end{equation}
但这正是单个 Bethe 根 $\sW=w$ 情形下的 Yang-Yang 函数 (\ref{yy})，只差一些仅依赖于 $z_a$ 而不依赖于 $\sW$ 的项；目前的推导对这些项并不敏感。

% [Tutor Note: 物理图像] 积掉 $\sY$ 的过程就是代数地解 $\partial\W/\partial\sY=0$ 再代回，在超势层面等价于在路径积分中对 $\sY$ 作鞍点近似；被积掉的 $\sY$ 的量子涨落给出与 $\zeta$ 无关的常数，故 (\ref{twig}) 中未出现它们。这一步也说明：$\zeta$ 起到了“把 $\log K$ 型势能引入”的作用，这正是 Yang-Yang 函数中 $\zeta\sum\log(\sW-z_a)$ 项的来源。

此时，读者想必希望看到通向含任意多个 Bethe 根的一般 Yang-Yang 函数的类似推导。对于有 $q$ 个光滑磁单极子的一般情形，我们用超场 $\sW_i$、$\sY_i$，$i=1,\dots,q$ 描述第 $i^{th}$ 个光滑磁单极子的位置。这些场的定义颇为微妙，将在第~\ref{scattering}~节讨论。$\sW_i$ 的期望值将成为第~\ref{opbethe}~节的 Bethe 根 $w_i$。与那里的讨论一样，引入多项式 $Q(z)=\prod_{i=1}^q(z-\sW_i)$ 是方便的。结果表明，把超势 (\ref{bigsuper}) 推广到任意多个光滑磁单极子给出
\begin{equation}\label{giantsuper}\W=\ca\sum_i \sW_i +\zeta\sum_i \sY_i+\Lambda
\sum_i \frac{K(W_i)}{Q'(W_i)}\exp(-Y_i).   \end{equation}

为还原第~\ref{teqone}~节的定性结果，我们首先取 $\zeta=0$、$\ca\not=0$。要找到临界点，各个 $\sW_i$ 必须等于 $K$ 的零点 $z_a$。假定所有磁荷 $k_a$ 均为 1，则至多只有一个 $\sW_i$ 可以等于同一个 $z_a$（否则 $Q'(\sW_i)$ 的零点会抵消 $K(\sW_i)$ 的零点，条件 $\partial\W/\partial \sY_i=0$ 就不成立）。对所有 $q$ 的取值求和，总共有 $2^d$ 个解——每个 $z_a$ 可以与某个 $\sW_i$ 相等，也可以不相等。若取 $\ca\to 0$，则除一个解（完全没有光滑磁单极子的那个）之外，其余解都以 $\sY_i\sim \log(1/\ca)$ 的方式消失。

对于 $\zeta\not=0$，正如推导 (\ref{twig}) 时那样，方便的做法是积掉大质量场 $\sY_i$。模去不依赖于 $\sW_i$ 的项，我们得到的超势恰是 Yang-Yang 函数：
\begin{equation}\label{longyang} \W=\sum_i\ca\sW_i+\zeta\sum_i\log(K(\sW_i)/Q'(\sW_i))
=\sum_i\ca\sW_i+\zeta\sum_{i,a}\log(\sW_i-z_a)-2\zeta\sum_{i<j}\log(\sW_i-\sW_j).\end{equation}

细心的读者可能注意到我们所述内容中的一处遗漏。在 $\zeta=0$ 的情形，我们没有分析某个 $a$ 满足 $k_a>1$ 的问题。为此，重要的是考虑对某对 $i,j$ 有 $W_i=W_j=z_a$ 的情形，但我们用来描述磁单极子模空间的这些坐标在 $W_i=W_j$ 时其实是不适用的。我们马上给出更好的描述。

% [Tutor Note: 易错点] 由 (\ref{giantsuper}) 出发，$\partial\W/\partial\sY_i=0$ 直接给出 $\exp(-\sY_i)$，但代入后 $\log K(\sW_i)-\log Q'(\sW_i)$ 的化简用到 $Q'(\sW_i)=\prod_{j\not=i}(\sW_i-\sW_j)$，从而产生 $-\sum_{j\not=i}\log(\sW_i-\sW_j)$，对称化后即 $-2\zeta\sum_{i<j}\log(\sW_i-\sW_j)$。系数 $2$ 来自每对 $(i,j)$ 被计算两次，切勿漏掉。

\subsection{磁单极子的坐标}\label{scattering}

$\R^3$ 上多个光滑 BPS 磁单极子的模空间是一个内容丰富的数学理论的对象~\cite{AH}。我们只需要其中很小一部分：当磁单极子在空间上彼此远离（或等价地，当对称破缺很强）时，模空间可以有有效的阿贝尔描述。在这一描述中，磁单极子被视为通过阿贝尔规范多重态彼此相互作用的“点状”Dirac 磁单极子。每个磁单极子在 $\R^3$ 中有一个位置，还有一个作为电荷转动集体坐标的角度变量 $\vartheta$。如第~\ref{outline}~节那样，一旦在模空间上选定某个复结构，每个磁单极子的位置与角度坐标就组合成一对手征超场 $\sW$、$\nY$。然而，$\nY$ 的定义有一处微妙之处，这也是我们把记号从第~\ref{outline}~节改换过来的原因。



坐标 $\nY$ 的角度部分参数化了在把光滑磁单极子解“粘合”到阿贝尔解之前对其作 $U(1)$ 规范变换的自由度。因此 $e^{\nY}$ 是 $\sW$ 平面上 $U(1)$ 规范丛的全纯截面。这就是为什么在边界上存在奇异磁单极子时，单个光滑磁单极子的指数型超势 $\exp(-\nY)$ 需要一个前因子 $K(\sW)$：超势应当是一个函数，而奇异磁单极子把 $e^{-\nY}$ 变成了丛 $\otimes_a {\cal O}(z_a)^{-k_a}$ 的截面，其指数是奇异磁单极子在阿贝尔有效场论中的磁荷。于是我们乘以 $K(\sW)$ 来补偿。更直观的解释是：超势编码了 BPS 磁单极子与复希格斯场 $\varphi$ 的非对角部分的相互作用，而 $\varphi$ 在 $z_a$ 处有 $k_a$ 阶零点。因此超势获得一个因子 $(z-z_a)^{k_a}$。

另一方面，在阿贝尔理论中，BPS 磁单极子表现得像一个奇异磁单极子。在任意给定的 $y$ 值处，可以把低能描述的 $U(1)$ 规范丛限制到 $\sW$ 平面上。当 $y$ 增大、越过某个 BPS 磁单极子所在位置时，$\sW$ 平面上的 $U(1)$ 规范丛发生跳变。换言之，位于第一个 BPS 磁单极子右侧的第二个 BPS 磁单极子会感受到第一个磁单极子所在位置 $\sW_1$ 处存在磁荷 $+2$ 的奇异磁单极子。因此我们预期超势为
\begin{equation}
\W=K(\sW_1)\exp(-\nY_1) + \frac{K(\sW_2)}{(\sW_2 - \sW_1)^2}\exp(-\nY_2)
\end{equation}
类似地，对于实部递增的多个磁单极子：
\begin{equation}
\W=\sum_i \frac{K(\sW_i)}{\prod_{j<i} (\sW_i - \sW_j)^2}\exp(-\nY_i)
\end{equation}

只要各 $\sW_i$ 互不相同，这个超势就与 (\ref{giantsuper}) 的相关部分等价。两者由如下变量替换相联系
\begin{equation} \label{rela}
\exp(-\nY_i) = \frac{\prod_{j<i} (\sW_i - \sW_j)}{\prod_{j>i} (\sW_i - \sW_j)} \exp(-\sY_i)
\end{equation}
这一重新定义不影响超势中关于 $\ca$ 与 $\zeta$ 线性的部分，因为 $\sum_i \nY_i = \sum_i \sY_i$。然而，实部才是磁单极子在 $y$ 方向实际位置的是 $\nY_i$ 而非 $\sY_i$；此外，当两个或更多 $\sW_i$ 重合时，$\nY_i$ 与 $\sY_i$ 之差发散。用 $\nY_i$ 表达的超势能正确重现 $\zeta =0$ 时对任意 $k_a$ 值解的计数。若 $\ca=0$、$K(z) = z^{k_a}$，超势为
\begin{equation}
\W=\sum_{i=1}^q \frac{\sW_i^{k_a}}{\prod_{j<i} (\sW_i - \sW_j)^2}\exp(-\nY_i)
\end{equation}
当所有 $\sW_i$ 为零且光滑磁单极子的数目 $q$ 不超过 $k/2$ 时，该函数对任意的 $\nY_i$ 都取极值，因为 $\sW_i$ 全为零时各前因子至少有 $2$ 阶零点。这重现了我们在第~\ref{solutions}~节所得的结果。若我们用 $\sY_i$ 作同样的计算，似乎即使在原点处磁单极子数目大于 $k_a/2$ 时也会得到解，但 $\sW_i$ 重合时 $\sY_i$ 不是好变量。$\ca\not=0,\,\zeta=0$ 的情形留给读者。对于 $\zeta\not=0$，$\sY_i$ 与 $\nY_i$ 之间的差别并不重要。

% [Tutor Note: 约定差异] 式 (\ref{rela}) 中的指数是 $\prod_{j<i}/\prod_{j>i}$ 之比，符号与顺序都不可随意改动：它来自散射矩阵乘积 $S=S_q\cdots S_1$ 的分解，$\nY_i$ 的实部才是第 $i$ 个磁单极子的 $y$ 坐标（按 $y$ 递减或递增排序）。若把比值写成 $\prod_{j<i}/\prod_{j>i}$ 的倒数，则 (\ref{rela}) 之后“$\sum_i\nY_i=\sum_i\sY_i$”的论断仍成立，但 $y$ 位置的解释会颠倒。

把低能描述中的坐标与磁单极子模空间的精确非阿贝尔描述相匹配是很有意思的。精确的磁单极子模空间由算符 $\D_y=D_y+i[\phi_t,\,\cdot\,]$ 的散射矩阵参数化~\cite{Hurtubise}。散射矩阵取如下形式
\begin{equation}
S = \begin{pmatrix} Q(z) & P(z) \cr \tilde P(z) & R(z)\end{pmatrix} \qquad QR - P \tilde P =1
\end{equation}
其中 $Q$ 是 $q$ 阶首一多项式，$P$ 与 $\tilde P$ 是次数至多为 $q-1$ 的多项式，$R$ 是次数至多为 $q-2$ 的多项式。$P$ 与 $Q$ 必然互素（这由条件 $QR-P\tilde P=1$ 推出），并且 $\tilde P$ 与 $R$ 都由 $Q$ 与 $P$ 唯一确定。可以通过 $Q(\sW_i)=0$ 与 $P(\sW_i) = \exp \sY_i$ 在磁单极子模空间上定义坐标 $\sW_i$、$\sY_i$。然而，当各 $\sW_i$ 不互不相同时，这个定义并不好用。

% [Tutor Note: 物理图像] 散射矩阵的四个方块中，$Q$ 的零点给出磁单极子位置 $\sW_i$，$P$ 在零点处的取值给出角度/集体坐标 $\sY_i$；条件 $QR-P\tilde P=1$ 即 Wronskian 型归一化，保证 $\det S=1$。这个“精确”描述与低能阿贝尔描述的关系正是 (\ref{rela})。

为找到在对称破缺很强时仍更好用的参数化，我们可以如下进行。对单个磁单极子，散射矩阵取如下形式
\begin{equation}
S_1 = \begin{pmatrix} (z - \sW_1) &~ e^{\nY_1} \cr -e^{- \nY_1} &0\end{pmatrix}
\end{equation}
这个表达式可以自然地与如下事实相匹配：在低能下，光滑磁单极子可以近似为奇异的 Dirac 磁单极子。它告诉我们，方程 $\D_y\psi=0$ 有两个解 $\psi_\pm$，其行为分别为
\begin{align}
\psi_+ &\sim e^{\ca_1 y/2}\begin{pmatrix}1 \cr 0 \end{pmatrix} 
\qquad & y<<0 \cr  \psi_+ &\sim (z-\sW_1) e^{\ca_1 y/2}
\begin{pmatrix}1 \cr 0 \end{pmatrix} + e^{\nY_1 -\ca_1 y/2}
\begin{pmatrix}0 \cr 1 \end{pmatrix}\qquad &  y>>0
\end{align}
\begin{align}
\psi_- &\sim e^{-\ca_1 y/2}\begin{pmatrix}0 \cr 1 \end{pmatrix}\qquad  & y>>0
 \cr  \psi_- &\sim (z-\sW_1) e^{-\ca_1 y/2}\begin{pmatrix}0 \cr 1 \end{pmatrix} 
 + e^{-\nY_1 +\ca_1 y/2}\begin{pmatrix}1 \cr 0 \end{pmatrix} \qquad & y<<0
\end{align}
这里 $\psi_+$（$\psi_-$）是唯一一个在 $y<<0$（$y>>0$）时很小的解。阿贝尔理论中的奇异磁单极子解会在 $y>>0$（$y<<0$）时给出相同的指数增长，而次指数修正来自光滑磁单极子解中指数衰减的修正项。

% [Tutor Note: 易错点] 上两式中 $y<<0$ 与 $y>>0$ 的渐近行为方向容易混淆：$\psi_+$ 在 $y\to-\infty$ 时按 $e^{\ca_1y/2}$ 衰减（当 $\ca_1>0$），在 $y\to+\infty$ 时含 $\sW_1$、$\nY_1$ 的项代表散射相移。注意 $\nY_1$ 以 $e^{+\nY_1}$ 出现在 $S_1$ 的右上角，却以 $e^{\nY_1-\ca_1y/2}$ 出现在 $\psi_+$ 的 $y>>0$ 行为中，核对时请以 $S$ 的定义为准。

对于多个彼此远离的光滑磁单极子的位形，散射矩阵是一个乘积
\begin{equation}\label{wellsep}
S = S_q S_{q-1} \cdots S_1.
\end{equation}
其中 $S_a$ 是第 $a^{th}$ 个磁单极子贡献的散射矩阵，各磁单极子按 $y$ 方向排序。用
\begin{equation}
S_a = \begin{pmatrix} (z - \sW_a) & ~e^{\nY_a} \cr -e^{- \nY_a} &0\end{pmatrix}.
\end{equation}
参数化模空间是自然的。这使我们能用 $\sW_i$ 与 $\nY_i$ 写出 $P$ 与 $Q$，最后利用 $P(\sW_i)=\exp(\sY_i)$ 把 $\sY_i$ 用 $\sW_i$ 与 $\nY_i$ 表达出来。

\subsection{通过 $M$ 理论与膜的实现}\label{realm}

这里我们将解释一种理解指数型超势 (\ref{forg}) 的 $M$ 理论途径。我们先作一些准备。
